% Native LaTeX handout. Edit teaching content and print layout in this file.
\documentclass[11pt,a4paper]{article}
\usepackage[margin=23mm]{geometry}
\usepackage{fontspec}
\usepackage{xeCJK}
\usepackage[indent=2em]{parskip}
\usepackage{setspace}
\usepackage{enumitem}
\usepackage[nobottomtitles*]{titlesec}
\usepackage{indentfirst}
\usepackage{needspace}
\usepackage{amsmath}
\usepackage{graphicx}
\usepackage{array}
\usepackage{booktabs}
\usepackage{tabularx}
\usepackage{longtable}
\usepackage{fvextra}
\usepackage{xcolor}
\usepackage{xurl}
\usepackage[unicode,hidelinks]{hyperref}

\setsansfont[BoldFont=HelveticaNeue-Bold,ItalicFont=HelveticaNeue-Italic]{HelveticaNeue}
\setmainfont[BoldFont=HelveticaNeue-Bold,ItalicFont=HelveticaNeue-Italic]{HelveticaNeue}
\setmonofont{Menlo-Regular}
\setCJKmainfont[BoldFont=IBMPlexSansTC-SmBld]{IBMPlexSansTC-Text}
\setCJKsansfont[BoldFont=IBMPlexSansTC-SmBld]{IBMPlexSansTC-Text}
\setCJKmonofont[BoldFont=IBMPlexSansTC-SmBld]{IBMPlexSansTC-Text}
\xeCJKsetup{AutoFallBack=true}
\setCJKfallbackfamilyfont{\CJKrmdefault}[BoldFont=IBMPlexSansSC-SmBld]{IBMPlexSansSC-Text}
\setCJKfallbackfamilyfont{\CJKsfdefault}[BoldFont=IBMPlexSansSC-SmBld]{IBMPlexSansSC-Text}
\setCJKfallbackfamilyfont{\CJKttdefault}[BoldFont=IBMPlexSansSC-SmBld]{IBMPlexSansSC-Text}
\definecolor{handoutaccent}{RGB}{89,43,31}
\setcounter{secnumdepth}{1}

% Increase the original paragraph leading by 25%: 1.12 * 1.25 = 1.40.
% Keep the existing list and code-block gaps independent of that leading.
\setstretch{1.4}
\setlist{leftmargin=*,topsep=6.8pt,itemsep=4.25pt,parsep=0pt}
\DefineVerbatimEnvironment{Code}{Verbatim}{
    fontsize=\small,baselinestretch=1.4,
    breaklines=true,breakanywhere=true,samepage=true,
    frame=single,rulecolor=\color{black!20},framesep=3mm}
% Add external space above code, including code directly inside a list item.
\BeforeBeginEnvironment{Code}{\par\addvspace{10.2pt}\nopagebreak[4]}
\DeclareRobustCommand{\code}[1]{{\small\ttfamily #1}}
\pdfstringdefDisableCommands{\def\code#1{#1}}
\newcommand{\AnswerLink}[1]{%
    \par\noindent\hyperref[#1]{\textcolor{handoutaccent}{查看參考答案（第 \pageref*{#1} 頁）}}\par}
\newcommand{\BackLink}[1]{%
    \par\noindent\hyperref[#1]{\textcolor{handoutaccent}{返回題目（第 \pageref*{#1} 頁）}}\par}
\newcommand{\PATImage}[1]{%
    \begin{center}
        % 60% of the previous display dimensions; preserve the original files.
        \includegraphics[width=0.552\linewidth,height=0.276\textheight,keepaspectratio]{img/#1}
    \end{center}}
\newenvironment{PATStep}{%
    \par\addvspace{0.7\baselineskip}\noindent
    \begin{minipage}{\linewidth}\setlength{\parindent}{2em}%
}{\end{minipage}\par}
\newcommand{\PATRow}[3]{%
    \par\noindent\begin{minipage}[t]{0.52\linewidth}
        \vspace{0pt}\setlength{\parindent}{2em}
        \subsubsection{#1}
        #2\par
    \end{minipage}\hfill
    \begin{minipage}[t]{0.43\linewidth}
        \vspace{0pt}\centering
        \includegraphics[width=\linewidth,height=0.25\textheight,keepaspectratio]{img/#3}
    \end{minipage}\par\addvspace{0.6\baselineskip}}
\titleformat{\section}{\Large\bfseries}{\thesection.}{0.6em}{}
\titleformat{\subsection}{\large\bfseries}{}{0pt}{}
\titleformat{\subsubsection}{\normalsize\bfseries}{}{0pt}{}
\titlespacing{\section}{0pt}{1.5\baselineskip}{0.6\baselineskip}
\titlespacing{\subsection}{0pt}{\baselineskip}{0.35\baselineskip}
\titlespacing{\subsubsection}{0pt}{0.7\baselineskip}{0.25\baselineskip}
\renewcommand{\bottomtitlespace}{0.15\textheight}
\setlength{\emergencystretch}{2em}
\hypersetup{pdftitle={Git Workshop Handout},pdfauthor={WHUAI Team},bookmarksnumbered=true}

\begin{document}

\begin{center}
    {\LARGE\bfseries\color{handoutaccent} Git 入門課：學生練習講義\par}
\end{center}

這份講義把課堂操作整理成可獨立閱讀、列印的版本。你可以搭配投影片使用，也可以依序閱讀；Sxx 是投影片對照 ID，方便回頭查找對應畫面。

五個練習使用同一個小專案，請先把教材倉庫中的 \code{exercises/grade-lab} 複製到教材倉庫外，當作獨立副本。每個練習都在該副本中進行；完成後再回到投影片或本講義核對概念。

\section{開始之前}

\subsection{課前準備}

\begin{itemize}
    \item 安裝 \href{https://git-scm.com/downloads}{Git}\footnote{\url{https://git-scm.com/downloads}}，版本 2.28 以上。Windows 使用 Git Bash；macOS / Linux 使用終端機。
    \item 準備可保存純文字的編輯器。使用 VS Code 且已安裝 \code{code} 命令時，可設定如下；已有其他可用編輯器就不必更換：
\begin{Code}
git config --global core.editor "code --wait"
\end{Code}
    \item 設定 commit 作者資訊時，把範例改成自己的姓名與電子郵件；這不是 GitHub 登入。也可使用 GitHub 設定頁提供的 noreply 電子郵件。
\begin{Code}
git --version
git config --global user.name "Alex Chen"
git config --global user.email "alex@example.com"
\end{Code}
    \item 準備可登入的 GitHub 帳號，確認網路能存取 GitHub。\hyperref[guide:pat]{第 4 章的 PAT 教學}會帶著完成第一次推送的認證設定；已有可透過 HTTPS 正常 \code{push} 的認證可繼續使用。
    \item Python 是選用的檢查工具，不是完成 Git 練習的必要條件。依環境使用 \code{python}、\code{python3} 或 \code{py}。
\end{itemize}

\subsection{Git 與 GitHub 的分工}

Git 是安裝在本機的版本控制工具；GitHub 是存放 Git 倉庫的遠端服務。沒有 GitHub 仍可 \code{init}、\code{add}、\code{commit}，只是 commit 只留在本機。GitHub 讓他人取得、審閱與備份遠端倉庫。

\subsection{準備獨立的練習目錄（S05）}

{\sloppy
用檔案管理器把本教材的 \code{exercises/grade-lab} \textbf{整個複製}到個人主目錄，改名為 \code{git-grade-lab}。目標必須是新目錄，不可覆蓋既有專案。已有同名目錄時，改用新名稱，並相應調整後續本機路徑。\par
}

打開 Git Bash / 終端機，進入副本：

\begin{Code}
cd "$HOME/git-grade-lab"
pwd
\end{Code}

應看到以下兩個起始檔案，沒有 \code{.git}：

\begin{Code}
git-grade-lab/
  README.md
  grades.py
\end{Code}

起始 \code{README.md}：

\begin{Code}
# Grade calculator

Calculate the average of a list of scores.
\end{Code}

起始 \code{grades.py}：

\begin{Code}
def average(scores):
    return sum(scores) / len(scores)


print("Average:", average([70, 80, 90]))
\end{Code}

可以向 AI 求助，但課堂上請親手輸入每一條命令。

每次編輯後先在編輯器保存。不要輸入 shell 的 \code{\$} 提示符號。除非另外說明，命令都在目前的練習倉庫中執行。

\section{保存與檢查版本}

\subsection{工作目錄與暫存區}

工作目錄是正在編輯的檔案；暫存區記錄下一次 commit 要包含的內容；本機倉庫保存已完成的 commit。\code{git add} 把工作目錄的變更放入暫存區，\code{git commit} 把暫存區寫成新 commit。

\code{git status} 顯示目前 branch、已暫存與未暫存變更、未追蹤檔案。\code{git diff} 比較工作目錄與暫存區；\code{git diff --staged} 比較暫存區與最後一次 commit（\code{HEAD}）。因此同一檔案可同時有這兩種差異。

\begin{center}
工作目錄 $\xrightarrow{\texttt{add}}$ 暫存區 $\xrightarrow{\texttt{commit}}$ 本機倉庫 $\xrightarrow{\texttt{push}}$ 遠端倉庫
\end{center}

\subsection{練習 1：分開保存兩次修改（S07–S11）}
\label{exercise:staging}

\subsubsection{建立第一個版本}

\begin{Code}
git init -b main
git status
git add README.md grades.py
git status
git commit -m "Add grade calculator"
git status
\end{Code}

預期目前位於 \code{main}，已有第一個 commit，工作目錄乾淨。

\subsubsection{先預測，再 commit}

\begin{enumerate}
    \item 在 README 末尾加入 \code{Input: a list of scores.}，保存。
    \item 執行 \code{git add README.md}。
    \item 再追加 \code{Output: the average score.}，保存，\textbf{先不要 \code{add}}。
    \item 執行下面三個命令，並預測下一個 commit 會包含哪一行。
\end{enumerate}

\begin{Code}
git status
git diff
git diff --staged
\end{Code}

先只 commit 輸入說明，再把輸出說明保存成另一個 commit。

完成條件：至少有三個 commit；Input 與 Output 分別 commit；工作目錄乾淨；能說明兩種 \code{diff} 各自比較的對象。

\AnswerLink{answer:staging}

\subsubsection{小檢查：忽略不需要分享的檔案（S12）}

建立 \code{.gitignore}，內容如下。註解必須獨立成行。

\begin{Code}
# Local credentials
.env
# Generated Python files
__pycache__/
*.pyc
\end{Code}

建立 \code{.env}，只放假資料 \code{DEMO\_ONLY=example}。確認 \code{.env} 沒有出現在一般 \code{status} 中，再保存共用規則：

\begin{Code}
git status
git add .gitignore
git commit -m "Ignore local and generated files"
git status
\end{Code}

這些規則對未追蹤檔案生效；已追蹤的檔案不會自動停止追蹤。忽略也不會清除歷史中的內容。

\subsubsection{一起暫存目前目錄的變更（S12a）}

當目前目錄與子目錄中的變更都要納入同一個 commit 時，可用 \code{git add .} 一起暫存，包括新增、修改與刪除。\code{.} 表示目前目錄；從專案根目錄執行時，會處理整個專案。

\begin{Code}
git status
git add .
git diff --staged
\end{Code}

忽略規則只會略過符合規則的未追蹤檔案；已追蹤的檔案仍會暫存。commit 前先讀 \code{git diff --staged}，確認暫存內容都應包含在本次 commit 中。只想選某幾個檔案時，仍可用 \code{git add 檔名}。

\subsection{查歷史與選復原方式}

\code{git log --oneline} 看 commit 順序，\code{git show HEAD} 看目前 commit 的內容。復原前先判斷變更在哪裡：尚未暫存、已暫存但未 commit，或已經 commit。三種情況用不同命令，不能互相取代。

取消暫存用 \code{git restore --staged README.md}，只把暫存區回復到 \code{HEAD}，工作目錄的修改仍在。丟棄未暫存修改用 \code{git restore README.md}，會用暫存區版本覆蓋工作目錄；這會失去修改，先確認不要了再執行。兩者是選項，不要依序全部執行。

已 commit 的變更若已分享，通常用 \code{git revert} 新增 revert commit 來抵銷，保留歷史。它不會刪掉舊 commit；若只是本機尚未分享的錯誤，才在理解後考慮其他整理方式。

\subsection{練習 2：撤銷一次錯誤，保留歷史（S13–S15）}

先確定工作目錄乾淨，把 README 的\textbf{第一行}改為 \code{\# Broken calculator}，其餘行保留。

\begin{Code}
git add README.md
git commit -m "Change heading by mistake"
git log --oneline
git show HEAD
\end{Code}

確認 \code{HEAD} 是剛才的錯誤 commit，再執行：

\begin{Code}
git revert --no-edit HEAD
git log --oneline
git status
\end{Code}

完成條件：標題恢復為 \code{\# Grade calculator}，輸入／輸出說明仍在，工作目錄乾淨；歷史中同時有錯誤 commit 與抵銷 commit。\code{revert} 會新增 revert commit，不會刪掉舊 commit。

以下兩種常見情境是\textbf{選項，不要接著依序執行}：

\noindent\begin{tabularx}{\linewidth}{@{}p{0.2\linewidth}p{0.43\linewidth}X@{}}
\toprule
想做的事 & 命令 & 工作內容 \\
\midrule
取消 README 的暫存 & \code{git restore --staged README.md} & 保留工作目錄修改 \\
\addlinespace
丟棄 README 未暫存的修改 & \code{git restore README.md} & 用暫存區版本覆蓋；先確定修改不要了 \\
\bottomrule
\end{tabularx}

\section{branch 與合併}

\subsection{branch、HEAD 與合併方向}

branch 名稱是指向某個 commit 的可移動標籤；建立 branch 不會複製整個目錄。\code{HEAD} 通常指向目前所在 branch，commit 後該 branch 前進，\code{HEAD} 也跟著指向新 commit。

\code{git switch -c handle-empty} 從目前 commit 建立 branch 並切過去；往後 commit 記在 \code{handle-empty}。\code{git merge handle-empty} 要在 \code{main} 上執行，因為合併目標是目前 branch，\code{handle-empty} 是來源。

本例 \code{main} 自 branch 建立後沒有新 commit，因此這次合併可以快轉：\code{main} 直接移動到 \code{handle-empty} 的 commit，不產生 merge commit。

\subsection{練習 3：在 branch 上開發，再合併回 \code{main}（S16–S19）}

以下操作只需執行一次。若已建立 \code{handle-empty} branch，直接核對完成條件，不要重複建立同名 branch。

\begin{Code}
git status
git switch -c handle-empty
\end{Code}

把 \code{grades.py} 的函式改為以下版本，保留原本的 \code{print}：

\begin{Code}
def average(scores):
    if not scores:
        return 0
    return sum(scores) / len(scores)
\end{Code}

今天的專案約定：空清單的平均值回傳 0。

\begin{Code}
git add grades.py
git commit -m "Handle empty score lists"
git switch main
\end{Code}

先開啟 \code{grades.py}：\code{main} 應還沒有空清單判斷。再合併：

\begin{Code}
git merge handle-empty
git branch
git log --oneline --graph --all
git status
\end{Code}

完成條件：目前在 \code{main}，檔案已有空清單判斷，工作目錄乾淨。\code{main} 與 \code{handle-empty} 此時指向同一 commit，來源 branch 名稱仍存在。

可選驗證；依環境替換 \code{python} 命令：

\begin{Code}
python -c "from grades import average; assert average([70,80,90]) == 80; assert average([]) == 0"
\end{Code}

\subsection{內容衝突與解決步驟}

合併時若兩邊修改同一段內容，Git 無法判斷保留哪一版，會停下來標示衝突。這是需要人工判斷的情況：先讀兩邊的用意，再寫出同時保留雙方資訊的版本，而不是只選一邊。

衝突標記中的 \code{HEAD} 是目前 branch（\code{main}）的版本，另一段是來源 branch（\code{improve-title}）的版本。編輯檔案時移除 \code{<<<<<<<}、\code{=======}、\code{>>>>>>>}，留下整理後的內容。

解決後執行 \code{git add README.md}，表示衝突已處理並放入暫存區；再執行 \code{git commit} 完成 merge commit。\code{add} 不會自動判斷內容好壞，\code{commit} 才建立合併結果。

\subsection{練習 4：整合兩邊的標題修改（S20–S23）}

開始前位於 \code{main}，工作目錄乾淨。若已建立 \code{improve-title} 並完成兩邊 commit，直接核對完成條件，不必重跑本節。

\begin{Code}
git switch -c improve-title
\end{Code}

把 README 第一行改為 \code{\# Team grade calculator}，其他內容保留，再執行：

\begin{Code}
git add README.md
git commit -m "Describe team calculator"
git switch main
\end{Code}

把 \code{main} 的 README 第一行改為 \code{\# Class grade calculator}，其他內容保留：

\begin{Code}
git add README.md
git commit -m "Describe class calculator"
git merge improve-title
\end{Code}

\textbf{預期 \code{merge} 會停下來，並提示 README 有衝突。} 執行 \code{git status}，打開 README，找到：

\begin{Code}
<<<<<<< HEAD
# Class grade calculator
=======
# Team grade calculator
>>>>>>> improve-title
\end{Code}

將整段替換成 \code{\# Team grade calculator for class}，保留其餘說明。確認沒有留下衝突標記，再執行：

\begin{Code}
git add README.md
git commit -m "Combine team and class descriptions"
git log --oneline --graph --all
git status
\end{Code}

完成條件：\code{main} 的標題保留 Team 與 class 的意思；輸入／輸出說明仍在；沒有衝突標記；有 merge commit 且工作目錄乾淨。

若想取消\textbf{尚未完成}的合併，可以執行 \code{git merge --abort}。本練習從乾淨工作目錄開始；不要假設它能恢復任何未保存的內容。

\section{分享與協作}

\subsection{遠端、推送與 PR}

\code{origin} 是 clone 時預設使用的遠端別名，也是常見習慣，不是固定必須使用的名稱；\code{main} 是 branch 名稱；\code{origin/main} 是本機記住的遠端 \code{main} 狀態。\code{commit} 只更新本機倉庫；\code{push} 把本機 commit 送到遠端並更新遠端 branch；\code{fetch} 會取得本機缺少的 commit 等物件，並更新 \code{origin/main} 等遠端追蹤資訊，不改工作目錄。

PR 是 GitHub 上的審閱與合併請求，不是 Git 命令。推送來源 branch 後開 PR；之後在同一個來源 branch commit 並 \code{push}，原 PR 會更新。PR 合併後，本機 \code{main} 還要用 \code{pull} 取得遠端結果。

遠端與本機各有新 commit 時，推送可能因非快轉被拒絕。這通常只是歷史分歧：先 \code{fetch}，再在乾淨工作目錄合併 \code{origin/main}。若兩邊沒有改到同一段內容，整合通常可自動完成；若 Git 無法判斷保留哪一版，會停下來標示衝突，需照練習 4 人工解決。

\subsection{設定 GitHub 推送認證：PAT}
\label{guide:pat}

PAT（Personal access token）是在 GitHub 透過 HTTPS 推送時，代替密碼使用的憑證；它不會改動 \code{user.name}、\code{user.email} 或 commit 的作者資訊。若已能透過 HTTPS 正常 \code{push}，可以略過這段設定。

以下使用自己的 GitHub 帳號建立 fine-grained PAT。畫面中的帳號、頭像與日期都是示例，請依自己的帳號和課程日期操作。

\subsubsection{1--2. 開啟帳號設定}

\noindent\begin{minipage}[t]{0.47\linewidth}
    \vspace{0pt}
    \noindent\textbf{1. Settings}

    登入 GitHub 後，點右上角頭像，再點 \mbox{Settings}。

    \PATImage{github-pat-01-account-settings.png}
\end{minipage}\hfill
\begin{minipage}[t]{0.47\linewidth}
    \vspace{0pt}
    \noindent\textbf{2. Developer settings}

    在 Settings 左側選單捲到最下方，點 \mbox{Developer settings}。

    \PATImage{github-pat-02-developer-settings.png}
\end{minipage}\par

\begin{PATStep}
\subsubsection{3. 建立 fine-grained token}

展開 Personal access tokens，選 Fine-grained tokens，再按 Generate new token。

\PATImage{github-pat-03-generate-token.png}
\end{PATStep}

\begin{PATStep}
\subsubsection{4. 填寫名稱、擁有者與到期日}

Token name 可填 \code{Git workshop}；Description 選填。Resource owner 選自己的帳號；Expiration 選 7 或 30 天，確保涵蓋練習期間。圖中的到期日期僅為示例。

\PATImage{github-pat-04-token-details.png}
\end{PATStep}

\clearpage
\begin{PATStep}
\PATRow{5. 選擇可存取的倉庫}{%
    本課為了示範方便，與圖中一樣選 All repositories，授權會涵蓋自己目前與未來的倉庫。

    \textbf{實際專案或正式環境請選 Only select repositories，只勾選需要的倉庫。}Public repositories 這個選項只有唯讀權限，無法用來 \code{push}。
}{github-pat-05-repository-access.png}

\PATRow{6. 加入 Contents 權限}{%
    在 Permissions 的 Repositories 分類按 Add permissions，搜尋 Contents 並勾選；下一步再設定讀寫權限。
}{github-pat-06-add-contents-permission.png}

\PATRow{7. 設定讀寫權限並產生 PAT}{%
    將 Contents 設為 Read and write；Metadata 保留預設的 Read-only 即可，最後按 Generate token。

    產生後立即複製並妥善保存 PAT，完整內容只會顯示一次。不要把它寫進 README、remote URL 或命令中，也不要分享含有 PAT 的截圖。
}{github-pat-07-read-write-permission.png}
\end{PATStep}

設定步驟可對照 \href{https://docs.github.com/en/authentication/keeping-your-account-and-data-secure/managing-your-personal-access-tokens}{GitHub 官方 PAT 說明}。
\clearpage

\subsection{練習 5：推送、審閱並更新同一個 PR（S24–S29）}

\subsubsection{推送本機倉庫}

確認目前在 \code{main}。到 GitHub 建立自己擁有的空倉庫 \code{git-grade-lab}，\textbf{不新增 README、\code{.gitignore} 或授權檔}。只使用不含敏感資料的課堂專案；公開倉庫方便同學查看與留言，私有倉庫需事先安排存取權限。

{\sloppy
把 \code{YOUR-USER} 換成自己的 GitHub 帳號；若倉庫另有名稱，也修改網址中的 \code{git-grade-lab}。\par
}

若要照本課直接輸入 PAT 的方式操作，先在目前的練習倉庫執行以下設定；已有可用認證者可略過這行。

\begin{Code}
git config --local credential.helper ""
\end{Code}

這行只對本練習倉庫停用 \href{https://git-scm.com/docs/gitcredentials#_configuration_options}{credential helper}，不改全域設定，也不儲存 PAT。之後 \code{push} 需要認證時，須再次輸入帳號與 PAT。

接著設定 HTTPS remote，推送本機的 \code{main}：

\begin{Code}
git remote add origin https://github.com/YOUR-USER/git-grade-lab.git
git push -u origin main
\end{Code}

若出現下列提示，Username 填 GitHub 帳號，Password 貼上 PAT。這裡的帳號不是 \code{user.name}，Password 也不是 GitHub 網站的登入密碼。

\begin{Code}
Username for 'https://github.com': YOUR-USER
Password for 'https://YOUR-USER@github.com':
\end{Code}

終端機貼上 PAT 時通常不顯示字元或星號，貼好後按 Enter 即可。後續仍使用原本的 \code{git push}；需要認證時，輸入同一個尚未過期的 PAT。

到 GitHub 確認可以看到 README、\code{grades.py} 與 commit 紀錄，且不應包含 \code{.env}。

\subsubsection{認證失敗時先檢查}

核對 PAT 是否過期、Resource owner 是否為自己的帳號、Repository access 是否包含目標倉庫，以及 Contents 是否設為 Read and write。PAT 過期時，回到上述設定頁重新建立。

\subsubsection{建立文件 branch 與 PR}

\begin{Code}
git switch -c docs-usage
\end{Code}

在 README 追加用法，例如：

\begin{Code}
## Usage

Run the calculator with Python:

~~~bash
python grades.py
~~~

Expected output: Average: 80.0
\end{Code}

若使用 \code{python3} 或 \code{py}，請相應修改說明。

\Needspace{8\baselineskip}
\begin{Code}
git diff
git add README.md
git commit -m "Explain how to run the calculator"
git push -u origin docs-usage
\end{Code}

在 GitHub 新增 PR，base 選 \code{main}，compare 選 \code{docs-usage}。描述可使用：

\begin{Code}
目的：讓第一次取得這個專案的同學知道如何執行。
變更：新增執行命令與預期輸出。
檢查：核對命令與檔名；若環境有 Python，實際執行並記錄結果。
\end{Code}

\subsubsection{雙人審閱與更新}

\begin{enumerate}
    \item 各自提出自己的 PR，交換連結。審閱者檢查用法、檔名與預期結果，提出一個具體建議；不需要對方的 \code{push} 權限。
    \item 作者依意見修改本機 README，並維持在 \code{docs-usage} branch。
    \item 保存後 commit 並推送，\textbf{原 PR 會更新}：
\begin{Code}
git add README.md
git commit -m "Clarify usage after review"
git push
\end{Code}
    \item 在原 PR 回應意見，並重新檢查 Files changed。
    \item 由有權限的人在 GitHub 合併 PR；本課可選 Create a merge commit。
    \item 在本機更新 \code{main}：
\begin{Code}
git switch main
git pull --ff-only
git status
\end{Code}
\end{enumerate}

完成條件：PR 有目的、檢查說明與審閱紀錄；已依意見更新同一個 PR；PR 已合併；本機 \code{main} 有最終用法，工作目錄乾淨。

\subsubsection{補充：\code{clone} 與遠端更新（S26–S27，可選）}

這個小節是選用補充，可在完成五個練習後閱讀；以下在新目錄取得既有倉庫，不再執行 \code{init}：

\Needspace{7\baselineskip}
\begin{Code}
git clone https://github.com/YOUR-USER/git-grade-lab.git git-grade-copy
cd git-grade-copy
git status
\end{Code}

\code{origin/main} 是本機記住的遠端狀態。\code{fetch} 會更新它，不會直接改動工作檔案。如果 \code{main} 兩邊各有新 commit，先確定工作目錄乾淨，再處理非快轉推送被拒絕的情況：

\begin{Code}
git fetch origin
git merge --no-edit origin/main
\end{Code}

若有衝突，照練習 4 編輯、\code{add}、\code{commit}；無衝突就直接完成整合，再執行 \code{git push}。認證失敗或沒有寫入權限是另外的問題，不適用這個處理方式。完成補充後回到原本的練習目錄；第二份 \code{clone} 只用於觀察遠端更新，不用來完成 PR 練習。

% Keep the review questions together; their answers are at the end.
\Needspace{14\baselineskip}
\section{課後想一想（S30）}
\label{exercise:review}

請先嘗試作答並記下想法：

\begin{enumerate}
    \item \code{add} 後又改檔案，下次 \code{commit} 會保存哪一版？
    \item \code{commit} 成功後，GitHub 一定看得到嗎？
    \item 將功能 branch 合併進 \code{main} 前，要切到哪裡？
    \item 編輯完有衝突的檔案後，還需要做什麼？
    \item 如何更新同一個 PR？
\end{enumerate}

\AnswerLink{answer:review}

\section{排查補充}

\subsection{依現象排查}

先確認目前目錄、branch 與 \code{git status}，再依現象檢查。下表只列先檢查什麼；處理時避免使用會覆寫他人歷史或丟棄未保存內容的命令。

\begin{longtable}{@{}p{0.25\linewidth}p{\dimexpr0.75\linewidth-2\tabcolsep\relax}@{}}
\toprule
現象 & 先檢查 \\
\midrule
\endfirsthead
\toprule
現象 & 先檢查 \\
\midrule
\endhead
\bottomrule
\endfoot
命令輸出停在分頁器 & 畫面底部是否顯示 \code{:}；按 \code{q} 離開。必要時用 \code{git --no-pager log --oneline} 取得單次輸出。 \\
\addlinespace
同一檔案同時有已暫存與未暫存修改 & 分別看 \code{git diff --staged} 與 \code{git diff}；確認下一次 commit 要包含哪一版，再決定 \code{add} 或保留工作目錄修改。 \\
\addlinespace
建立 branch 時提示同名 branch 已存在 & 用 \code{git branch} 確認；切到既有 branch，或改用新名稱，不要重複建立。 \\
\addlinespace
推送被拒絕且提示非快轉 & 先 \code{git fetch}，檢查 \code{origin/main} 與本機 \code{main} 的差異；在工作目錄乾淨時合併，再推送。不要用強制推送。 \\
\addlinespace
認證失敗或沒有寫入權限 & 確認 remote URL 的協定（HTTPS 或 SSH）與認證方式相符；講義範例多用 HTTPS，但既有可用的 SSH 也可使用。確認帳號對該遠端倉庫有權限；不要把 Token 放進網址。 \\
\end{longtable}

\subsection{需要幫忙時}

先說明目前目錄、branch、想做的事與錯誤訊息，並附上 \code{git status} 與相關 \code{diff}。可查文件，或問同學、助教、AI；執行建議前先理解它會改動哪個狀態，不要直接貼出憑證或私人內容。

{\small
\noindent\begin{tabularx}{\linewidth}{@{}l>{\raggedright\arraybackslash}X@{}}
    \href{https://git-scm.com/book/zh-tw/v2}{Pro Git 繁體中文} & \url{https://git-scm.com/book/zh-tw/v2} \\
    \href{https://git-scm.com/docs/git-diff}{git diff} & \url{https://git-scm.com/docs/git-diff} \\
    \href{https://git-scm.com/docs/git-restore}{git restore} & \url{https://git-scm.com/docs/git-restore} \\
    \href{https://git-scm.com/docs/git-revert}{git revert} & \url{https://git-scm.com/docs/git-revert} \\
    \href{https://git-scm.com/docs/gitignore}{gitignore} & \url{https://git-scm.com/docs/gitignore} \\
    \href{https://git-scm.com/docs/git-pull}{git pull} & \url{https://git-scm.com/docs/git-pull} \\
    \href{https://docs.github.com/en/pull-requests}{GitHub Pull Requests} & \url{https://docs.github.com/en/pull-requests} \\
\end{tabularx}
}

\clearpage
\section{參考答案}

\subsection{練習 1：分開保存兩次修改（S07–S11）}
\label{answer:staging}

完成預測後再讀以下操作。

\code{git diff} 顯示未暫存的 Output；\code{git diff --staged} 顯示已暫存的 Input。同一個檔案可以同時有這兩種變更。

\begin{Code}
git commit -m "Document calculator input"
git diff
git add README.md
git commit -m "Document calculator output"
git log --oneline
git status
\end{Code}

第一次 \code{commit} 後，Output 仍在工作目錄，並未包含在那次 commit 中。

\BackLink{exercise:staging}

\subsection{課後想一想（S30）}
\label{answer:review}

以下逐題對照。

\begin{enumerate}
    \item 暫存區中最後一次 \code{add} 當時的內容。
    \item 不一定，還需 \code{push} 到遠端。
    \item \code{main}；\code{merge} 的目標是目前所在 branch。
    \item 檢查內容、執行 \code{git add}，再執行 \code{git commit} 完成合併。
    \item 在 PR 的來源 branch 修改、commit 並 \code{push}，原 PR 會更新。
\end{enumerate}

\BackLink{exercise:review}

\end{document}
